% Л9 -- Эндпоинты, контракт и JSON. Slides.
%
% Companion to the spoken script in ../L09-endpoints-and-json.md, which is the
% source of truth: the slides carry diagrams, code, real terminal output and
% short marks, and the delivery lives in the script only -- this deck carries
% no speaker notes. Slide text is in Russian because that is the content;
% everything else in this file is English.
%
% Build:  make          (latexmk -pdf; engine: pdflatex)
%
% Every C# listing was compiled on dotnet SDK 10.0 in a scratch
% `dotnet new web -n HelloWorld`; terminal output on the slides is captured
% from runs of that project. See ../CLAUDE.md for the verification table.
%
% House rules: ../CLAUDE.md and ../../../CLAUDE.md. Theme: ../../../../theme/.

\documentclass[aspectratio=169,10pt,t]{beamer}

\usetheme{dotnet}
\usetikzlibrary{arrows.meta,positioning,calc,decorations.pathreplacing,fit}
\usepackage{array}   % ragged-right table columns

% A ragged-right paragraph column: justified cells in narrow columns hyphenate
% badly, and a split word is the one thing the back row cannot read.
\newcolumntype{L}[1]{>{\raggedright\arraybackslash}p{#1}}

\dotnetlecture{Л9}

\title[Л9. Эндпоинты, контракт и JSON]{Эндпоинты, контракт и JSON}
\author{Дмитрий Курочкин}
\date{}

% Code sizes, one step down from \footnotesize. Smaller than 6.6pt does not
% carry to the back of the room.
\newcommand{\codesm}{\ttfamily\fontsize{7.6}{9.0}\selectfont}
\newcommand{\codexs}{\ttfamily\fontsize{7.0}{8.3}\selectfont}
\newcommand{\codexxs}{\ttfamily\fontsize{6.8}{7.8}\selectfont}
\newcommand{\termsm}{\ttfamily\fontsize{6.6}{7.9}\selectfont}

% A mark at the start of a listing line, for labels set beside the listing.
% Used through the listing's escapechar; the labels are drawn by an overlay
% picture after the listing, so the build needs two passes (latexmk does that).
\newcommand{\mk}[1]{\tikz[remember picture,overlay]\coordinate (#1) at (0,0.55ex);}

\tikzset{
  rawbrace/.style={decorate,decoration={brace,amplitude=4pt},hairline,
                   line width=0.8pt},
  rawlab/.style={anchor=west,font=\footnotesize,inner sep=0pt},
  num/.style={anchor=west,font=\small\bfseries,text=accent,inner sep=0pt},
  box/.style={draw=hairline,line width=0.8pt,fill=codebg,inner sep=4pt,
              align=center,minimum height=0.9cm},
  arr/.style={-{Straight Barb[length=1.6mm]},muted,line width=0.7pt},
  arra/.style={-{Straight Barb[length=1.8mm]},accent,line width=0.9pt},
  arrl/.style={-{Straight Barb[length=1.8mm]},alt,line width=0.9pt},
}

\begin{document}
\shorthandoff{"}

% ------------------------------------------------------------------ title --
% 1
\begin{frame}[plain]
  \titlepagednc{Лекция 9}{Эндпоинты, контракт и JSON}{}{Дмитрий Курочкин}
\end{frame}

% ============================================== 1. TWO WAYS TO DESCRIBE ==
\section{1. Два способа описать эндпоинт}

% 2
\begin{frame}[plain]
  \sectionopener{1}{Два способа описать эндпоинт}{Обработчик --- это
    \hi{делегат}.}
\end{frame}

% 3
\begin{frame}[plain]
  \vote{Кто видел классы с атрибутами над методами, словом
    \texttt{Controller} в имени и базовым классом?}
\end{frame}

% 4
\begin{frame}[fragile]{Минимальный API}
  \vspace{0.3em}
\begin{lstlisting}[style=csharp,basicstyle=\termsm,escapechar=|]
|\mk{a1}|var orders = app.MapGroup("/orders");

|\mk{a2}|orders.MapGet("/{id:int}", async Task<Results<Ok<OrderDto>, NotFound>> (
|\mk{a3}|    int id, IOrderService service, CancellationToken ct) =>
{
    var order = await service.FindAsync(id, ct);
|\mk{a4}|    return order is null
        ? TypedResults.NotFound()
        : TypedResults.Ok(order);
});

orders.MapPost("/", async Task<Created<OrderDto>> (
    CreateOrderDto dto, IOrderService service, CancellationToken ct) =>
{
    var created = await service.CreateAsync(dto, ct);
    return TypedResults.Created($"/orders/{created.Id}", created);
});
\end{lstlisting}
  \begin{tikzpicture}[remember picture,overlay,every node/.style={rawlab}]
    \node at ([xshift=10.7cm]a1) {группа};
    \node[text=accent] at ([xshift=10.7cm]a2) {результат в типе};
    \node at ([xshift=10.7cm]a3) {зависимости};
    \node at ([xshift=10.7cm]a4) {готовый ответ};
  \end{tikzpicture}

  \aside{Всё об обработчике --- в \hi{одной} сигнатуре.}
\end{frame}

% 5
\begin{frame}[fragile]{Группы}
  \slidetop
  \claim{Прицеплено к внешней --- \hi{наследуют} внутренние.}

\begin{lstlisting}[style=csharp,basicstyle=\codesm,escapechar=|]
|\mk{g1}|var api = app.MapGroup("/api/v1").RequireAuthorization();

|\mk{g2}|var orders = api.MapGroup("/orders");
|\mk{g3}|var users = api.MapGroup("/users");

orders.MapGet("/stats", () => "public stats")
|\mk{g4}|      .AllowAnonymous();
\end{lstlisting}
  \begin{tikzpicture}[remember picture,overlay,every node/.style={rawlab}]
    \node[text=accent] at ([xshift=9.6cm]g1) {весь API закрыт одной строкой};
    \draw[rawbrace] ([xshift=9.35cm,yshift=0.35em]g2) -- ([xshift=9.35cm,yshift=-0.3em]g3);
    \node at ([xshift=9.6cm]$(g2)!0.5!(g3)$) {префикс и требования};
    \node at ([xshift=9.6cm]g4) {исключение --- на маршруте};
  \end{tikzpicture}

  \slidegap
  \aside{К группе цепляют общее: авторизацию, фильтры, теги документа.}
\end{frame}

% 6
\begin{frame}{Откуда параметры}
  \slidetop
  \claim{Источник --- по \hi{типу и имени}, без атрибутов.}

  {\renewcommand{\arraystretch}{1.5}
  \begin{tabular}{@{}L{4.4cm}L{4.6cm}L{3.8cm}@{}}
    \colhead{muted}{параметр} & \colhead{muted}{откуда} & \colhead{muted}{как понял} \\
    \hi{\texttt{int id}}               & путь               & имя как в \texttt{\{id\}} \\
    \hi{\texttt{IOrderService service}} & контейнер зависимостей & тип зарегистрирован \\
    \hi{\texttt{CancellationToken ct}} & \texttt{RequestAborted} & служебный тип \\
    \hi{\texttt{CreateOrderDto dto}}   & тело запроса, JSON & сложный тип \\
  \end{tabular}}

  \slidegap
  \aside{Атрибут нужен там, где угадать нельзя: заголовок, другое имя.}
\end{frame}

% 7
\begin{frame}[fragile]{Возвращаемый тип}
  \slidetop
  \claim{Статусы --- \hi{в сигнатуре}, а не в комментарии.}

\begin{lstlisting}[style=csharp,basicstyle=\ttfamily\small]
Task<Results<Ok<OrderDto>, NotFound>>
\end{lstlisting}

  \rulegap
  \versus{TypedResults}{%
    \begin{cbul}{accent}
      \item конкретные типы: \texttt{Ok<T>}, \texttt{NotFound}
      \item складываются в \texttt{Results<...>}
      \item статусы видят документ и компилятор
    \end{cbul}}{Results}{%
    \begin{cbul}{alt}
      \item общий \texttt{IResult}
      \item по сигнатуре не понять, что вернёт
      \item статусы --- только в атрибутах
    \end{cbul}}

  \rulegap
  \aside{Без результата: объект --- 200 с JSON, строка --- текст, \texttt{Task} ---
    пустой 200.}
\end{frame}

% 8
\begin{frame}[fragile]{Контроллер}
  \vspace{0.3em}
\begin{lstlisting}[style=csharp,basicstyle=\codexxs,escapechar=|]
|\mk{c1}|[ApiController]
|\mk{c2}|[Route("orders")]
public sealed class OrdersController(IOrderService service)
|\mk{c3}|    : ControllerBase
{
|\mk{c4}|    [HttpGet("{id:int}")]
    public async Task<ActionResult<OrderDto>> Get(int id, CancellationToken ct)
    {
        var order = await service.FindAsync(id, ct);
|\mk{c5}|        return order is null ? NotFound() : Ok(order);
    }

    [HttpPost]
    public async Task<ActionResult<OrderDto>> Create(
        CreateOrderDto dto, CancellationToken ct)
    {
        var created = await service.CreateAsync(dto, ct);
        return CreatedAtAction(nameof(Get), new { id = created.Id }, created);
    }
}
\end{lstlisting}
  \begin{tikzpicture}[remember picture,overlay,every node/.style={rawlab}]
    \node[text=accent] at ([xshift=10.6cm]c1) {включает соглашения};
    \node at ([xshift=10.6cm]c2) {маршрут на классе};
    \node at ([xshift=10.6cm]c3) {базовый класс};
    \node at ([xshift=10.6cm]c4) {метод HTTP и путь};
    \node at ([xshift=10.6cm]c5) {ответы базового класса};
  \end{tikzpicture}
\end{frame}

% 9
\begin{frame}{Соглашения}
  \slidetop
  \claim{Правила, которых \lo{нет в коде}.}

  {\renewcommand{\arraystretch}{1.5}
  \begin{tabular}{@{}L{4.6cm}L{8.2cm}@{}}
    имя на \hi{\texttt{Controller}}       & класс --- контроллер \\
    \hi{публичный} метод                 & обработчик \\
    сложный тип в параметре              & тело запроса \\
    \texttt{ActionResult<T>}              & вернуть можно что угодно --- статусы \lo{не
      видны} \\
    экземпляр на \hi{запрос}             & поле живёт один запрос: счётчик в нём всегда
      ноль \\
  \end{tabular}}

  \slidegap
  \aside{Атрибуты уточняют соглашения там, где они не подходят.}
\end{frame}

% 10
\begin{frame}[plain]
  \vote{Мы нигде не зарегистрировали контроллер. Откуда фреймворк узнал о нём
    и о его маршрутах?}
\end{frame}

% 11
\begin{frame}{Как находят контроллеры}
  \slidetop
  \claim{При старте, один раз, через \hi{рефлексию}.}

  \begin{center}
    \begin{tikzpicture}
      \foreach \i/\t in {1/сборка,2/метаданные,3/все типы,
                         4/{\texttt{ControllerBase}\\или имя},5/атрибуты\\методов} {
        \node[box,minimum width=1.95cm,minimum height=1.1cm,inner sep=1pt,font=\scriptsize]
          (r\i) at ({(\i-1)*2.32},0) {\t};
      }
      \node[box,draw=accent,line width=1.0pt,minimum width=1.95cm,minimum height=1.1cm,
            inner sep=1pt,font=\scriptsize\bfseries] (r6) at (11.6,0) {таблица\\маршрутов};
      \foreach \i [evaluate=\i as \j using int(\i+1)] in {1,...,5}
        \draw[arr] (r\i) -- (r\j);
    \end{tikzpicture}
  \end{center}

  \slidegap
  \begin{cbul}{accent}
    \item сборка описывает себя сама: каждый тип, каждый метод, каждый атрибут
    \item подходящий класс --- наследует \texttt{ControllerBase} или назван на
      \texttt{Controller}
    \item вы написали класс --- и он работает, \hi{ничего не регистрируя}
  \end{cbul}
\end{frame}

% 12
\begin{frame}{Цена}
  \slidetop
  \versus{Минимальный API}{%
    \begin{cbul}{accent}
      \item прямая ссылка на делегат
      \item генератор кода вместо рефлексии
      \item AOT и обрезка --- \hi{да}
      \item легче в бенчмарках
    \end{cbul}}{Контроллеры}{%
    \begin{cbul}{alt}
      \item находятся рефлексией по имени
      \item обрезка выкинет: ссылок на класс нет
      \item AOT --- \lo{не поддерживается}
      \item фильтры и активация на каждом запросе
    \end{cbul}}

  \slidegap
  \hrulethin{0.42}
  \rulegap
  \aside{Генератор видит лямбду в \texttt{Map} и метод по имени. Делегат в
    переменной в цикле --- откат на рефлексию.}
\end{frame}

% 13
\begin{frame}{Что берём}
  \slidetop
  \claim{Новый сервис --- \hi{минимальный API}.}

  {\renewcommand{\arraystretch}{1.6}
  \begin{tabular}{@{}L{4.2cm}L{8.6cm}@{}}
    \hi{\textbf{всё в одном месте}} & путь, параметры, зависимости, результат ---
      одна сигнатура \\
    \hi{\textbf{честность типов}}   & статусы проверяет компилятор и читает документ \\
    \hi{\textbf{AOT и обрезка}}     & меньше магии, даже если AOT не нужен \\
  \end{tabular}}

  \slidegap
  \aside{Контроллеры не устарели. Проект уже на них, нужны их механизмы, команда
    знает их --- переписывать не нужно.}
\end{frame}

% 14
\begin{frame}[fragile]{Обработчики по ресурсам}
  \vspace{0.3em}
\begin{lstlisting}[style=csharp,basicstyle=\codexxs,escapechar=|]
public static class OrderEndpoints
{
|\mk{o1}|    public static IEndpointRouteBuilder MapOrders(
        this IEndpointRouteBuilder app)
    {
|\mk{o2}|        var orders = app.MapGroup("/orders").WithTags("Orders");

|\mk{o3}|        orders.MapGet("/{id:int}", GetById).WithName("GetOrder");
        orders.MapPost("/", Create);
        return app;
    }

|\mk{o4}|    static async Task<Results<Ok<OrderDto>, NotFound>> GetById(
        int id, IOrderService service, CancellationToken ct) =>
        await service.FindAsync(id, ct) is { } order
            ? TypedResults.Ok(order)
            : TypedResults.NotFound();
}

|\mk{o5}|app.MapOrders();   // Program.cs
\end{lstlisting}
  \begin{tikzpicture}[remember picture,overlay,every node/.style={rawlab}]
    \node at ([xshift=9.9cm]o1) {расширение на ресурс};
    \node at ([xshift=9.9cm]o2) {метаданные цепочкой};
    \node[text=accent] at ([xshift=9.9cm]o3) {обработчик по имени};
    \node at ([xshift=9.9cm]o4) {зовётся из теста};
    \node at ([xshift=9.9cm]o5) {файл запуска --- список};
  \end{tikzpicture}
\end{frame}

% 15
\begin{frame}[plain]
  \keyline{Обработчик --- делегат, и вся его сигнатура на виду.\\
    Контроллеры дают то же через соглашения и рефлексию.}
\end{frame}

% ============================================== 2. FROM REQUEST TO TYPES ==
\section{2. От запроса к типам .NET}

% 16
\begin{frame}[plain]
  \sectionopener{2}{От запроса к типам .NET}{Три шага, и \hi{каждый} может
    сломаться.}
\end{frame}

% 17
\begin{frame}{Три шага}
  \slidetop
  \begin{center}
    \begin{tikzpicture}[font=\footnotesize,
      st/.style={box,minimum width=2.25cm,minimum height=1.2cm,inner sep=2pt}]
      \node[st] (s1) at (0,0)     {запрос\\{\scriptsize строки, байты}};
      \node[st] (s2) at (2.85,0)  {маршрутизация\\{\scriptsize найти обработчик}};
      \node[st] (s3) at (5.7,0)   {привязка\\{\scriptsize достать данные}};
      \node[st] (s4) at (8.55,0)  {разбор\\{\scriptsize строку в тип}};
      \node[st,draw=accent,line width=1.0pt] (s5) at (11.4,0)
        {параметры\\{\scriptsize обработчика}};
      \foreach \a/\b in {s1/s2,s2/s3,s3/s4,s4/s5} \draw[arra] (\a) -- (\b);
      \node[font=\scriptsize,text=alt,below=0.25cm,align=center] at (s2.south)
        {маршрут\\не совпал};
      \node[font=\scriptsize,text=alt,below=0.25cm,align=center] at (s3.south)
        {источник\\не тот};
      \node[font=\scriptsize,text=alt,below=0.25cm,align=center] at (s4.south)
        {разобрано\\не так};
    \end{tikzpicture}
  \end{center}

  \slidegap
  \aside{Всё, что пришло из пути, строки запроса и заголовков, --- строки. Тело ---
    байты. Параметр --- число, дата, перечисление, объект.}
\end{frame}

% 18
\begin{frame}{Шаблоны маршрутов}
  \slidetop
  \claim{Фигурные скобки --- \hi{параметр}, всё остальное --- буквально.}

  {\renewcommand{\arraystretch}{1.35}
  \begin{tabular}{@{}L{5.6cm}L{7.2cm}@{}}
    \texttt{/orders}                  & литерал: совпадает буквально \\
    \texttt{/orders/\{id\}}           & параметр: любой сегмент \\
    \hi{\texttt{/orders/\{id:int\}}}  & с ограничением: только целое \\
    \texttt{/orders/\{id:int:min(1)\}} & два ограничения \\
    \texttt{/files/\{*path\}}         & остаток пути, с косыми чертами \\
    \texttt{/orders/\{page=1\}}       & значение по умолчанию \\
  \end{tabular}}
\end{frame}

% 19
\begin{frame}[fragile]{Ограничение выбирает маршрут}
  \slidetop
  \claim{Ограничение --- это \hi{выбор маршрута}, а не валидация.}

\begin{lstlisting}[style=term,basicstyle=\termsm]
$ curl -i http://localhost:5090/orders/abc
HTTP/1.1 404 Not Found
Content-Length: 0
Date: Tue, 10 Nov 2026 11:10:57 GMT
Server: Kestrel
\end{lstlisting}

  \slidegap
  {\renewcommand{\arraystretch}{1.25}\small
  \begin{tabular}{@{}L{4.6cm}L{8.2cm}@{}}
    \hi{\texttt{/orders/latest}}       & литерал --- \hi{важнее} всех \\
    \texttt{/orders/\{id:int\}}        & параметр с ограничением \\
    \texttt{/orders/\{id\}}            & параметр без ограничения \\
  \end{tabular}}

  \rulegap
  \aside{Два равных шаблона --- \lo{\texttt{AmbiguousMatchException}} на запросе,
    не на старте. Анализатор предупредит при сборке, если видит оба рядом.}
\end{frame}

% 20
\begin{frame}[fragile]{Что маршрутизация делает сама}
  \slidetop
  {\renewcommand{\arraystretch}{1.3}
  \begin{tabular}{@{}L{4.6cm}L{8.2cm}@{}}
    регистр пути           & \texttt{/Orders/1} --- это \texttt{/orders/1} \\
    косая черта в конце     & \texttt{/orders/1/} --- это \texttt{/orders/1} \\
    имя в строке запроса    & \texttt{Page} --- это \texttt{page} \\
    \hi{метод не тот}       & \hi{405} и заголовок \texttt{Allow} \\
  \end{tabular}}

  \slidegap
\begin{lstlisting}[style=term,basicstyle=\termsm]
$ curl -i -X PUT http://localhost:5090/orders/1
HTTP/1.1 405 Method Not Allowed
Content-Length: 0
Date: Tue, 10 Nov 2026 11:10:57 GMT
Server: Kestrel
Allow: GET
\end{lstlisting}

  \rulegap
  \aside{\texttt{HEAD} на маршрут с \texttt{MapGet} --- тоже \lo{405}: его
    разрешают явно.}
\end{frame}

% 21
\begin{frame}[fragile]{Привязка}
  \vspace{0.3em}
\begin{lstlisting}[style=csharp,basicstyle=\codexs,escapechar=|]
app.MapGet("/orders/{id:int}/items", (
|\mk{p1}|    int id,
|\mk{p2}|    IOrderService service,
|\mk{p3}|    CancellationToken ct,
|\mk{p4}|    int page = 1,
|\mk{p5}|    [FromQuery(Name = "size")] int pageSize = 20,
|\mk{p6}|    [FromHeader(Name = "X-Request-Id")] string? requestId = null)
    => service.GetItemsAsync(id, page, pageSize, ct));

app.MapPost("/orders", (
|\mk{p7}|    CreateOrderDto dto, IOrderService service, CancellationToken ct)
    => service.CreateAsync(dto, ct));
\end{lstlisting}
  \begin{tikzpicture}[remember picture,overlay,every node/.style={rawlab}]
    \node at ([xshift=10.5cm]p1) {из пути};
    \node at ([xshift=10.5cm]p2) {из контейнера};
    \node at ([xshift=10.5cm]p3) {из контекста};
    \node at ([xshift=10.5cm]p4) {строка запроса};
    \node at ([xshift=10.5cm]p5) {под другим именем};
    \node at ([xshift=10.5cm]p6) {из заголовка};
    \node[text=accent] at ([xshift=10.5cm]p7) {\texttt{dto} --- из тела};
  \end{tikzpicture}

  \slidegap
  \aside{Сначала обязательные, потом со значениями по умолчанию --- как у любого
    метода. Значение по умолчанию делает параметр \hi{необязательным}.}
\end{frame}

% 22
\begin{frame}[fragile]{Источники}
  \slidetop
  {\renewcommand{\arraystretch}{1.22}\small
  \begin{tabular}{@{}L{3.4cm}L{9.4cm}@{}}
    \hi{путь}           & имя совпало с параметром шаблона \\
    \hi{строка запроса} & простой тип не из пути; \texttt{ids=1\&ids=2} --- массив \\
    заголовок           & только с \texttt{[FromHeader]}: по имени не угадать \\
    \hi{тело}           & сложный тип, JSON; такой параметр --- \lo{один} \\
    форма               & \texttt{[FromForm]}: HTML-форма, файлы \\
  \end{tabular}}

  \slidegap
\begin{lstlisting}[style=term,basicstyle=\termsm]
$ curl -i -d @order.json http://localhost:5090/orders
HTTP/1.1 415 Unsupported Media Type
Content-Length: 0
Date: Tue, 10 Nov 2026 11:10:57 GMT
Server: Kestrel
\end{lstlisting}

  \rulegap
  \aside{«Тело приходит пустым» --- сначала \hi{\texttt{Content-Type}}:
    \texttt{curl -d} шлёт форму. Десять параметров --- один объект с
    \texttt{[AsParameters]}.}
\end{frame}

% 23
\begin{frame}[fragile]{Свой тип в параметре}
  \vspace{0.3em}
\begin{lstlisting}[style=csharp,basicstyle=\codexs,escapechar=`]
public readonly record struct Email(string Value)
{
`\mk{e1}`    public static bool TryParse(string? s, out Email email)
    {
        email = default;
        if (string.IsNullOrWhiteSpace(s) || !s.Contains('@')) return false;
        email = new Email(s.Trim().ToLowerInvariant());
        return true;
    }
}

`\mk{e2}`app.MapGet("/customers/by-email", (Email email) => new CustomerDto(email, "Anna"));
\end{lstlisting}
  \begin{tikzpicture}[remember picture,overlay,every node/.style={rawlab}]
    \node[text=accent] at ([xshift=9.2cm]e1) {фреймворк найдёт сам};
  \end{tikzpicture}

\begin{lstlisting}[style=term,basicstyle=\termsm]
$ curl -s 'http://localhost:5090/customers/by-email?email=Anna@Example.com'
{"email":"anna@example.com","name":"Anna"}
\end{lstlisting}

  \rulegap
  \aside{Формат проверен при привязке, \hi{один раз}. \texttt{email=bad} --- 400.}
\end{frame}

% 24
\begin{frame}[plain]
  \vote{Параметр типа \texttt{decimal}, в строке запроса пришло
    \texttt{12,5}. Что получит обработчик?}
\end{frame}

% 25
\begin{frame}[fragile]{Четыре ловушки разбора}
  \slidetop
\begin{lstlisting}[style=term,basicstyle=\termsm]
$ curl -s 'http://localhost:5090/price?amount=12,5'
125
\end{lstlisting}

  \slidegap
  {\renewcommand{\arraystretch}{1.35}
  \begin{tabular}{@{}L{4.6cm}L{8.2cm}@{}}
    \texttt{decimal}: \texttt{12,5}    & \lo{125}: запятая --- разделитель тысяч \\
    \texttt{DateTime}: \texttt{03/04/2026} & \lo{4 марта}: месяц первым \\
    \texttt{enum}: \texttt{99}, \texttt{paid} & значение, которого нет; 400 \\
    \texttt{int page = 1}: \texttt{page=} & 400, а \lo{не} умолчание \\
  \end{tabular}}

  \slidegap
  \aside{Разбор по \hi{инвариантной} культуре --- одинаково на любой машине, и на
    машине с русской локалью тоже. Формат --- в контракте, всё вне формата ---
    отказ.}
\end{frame}

% 26
\begin{frame}{DTO $\neq$ сущность}
  \slidetop
  \begin{center}
    \begin{tikzpicture}[font=\footnotesize]
      \node[box,minimum width=4.6cm,minimum height=2.7cm,align=left] (e) at (0,0)
        {\ttfamily Id\\\ttfamily CustomerEmail\\\ttfamily Status\\\ttfamily Items\\
         \ttfamily\lo{CostPrice}\\\ttfamily\lo{InternalNote}};
      \node[box,draw=accent,line width=1.0pt,minimum width=4.6cm,minimum height=2.7cm,
            align=left] (d) at (8.2,0)
        {\ttfamily Id\\\ttfamily CustomerEmail\\\ttfamily Status\\\ttfamily Total};
      \node[font=\scriptsize\bfseries,text=muted,above=0.3em] at (e.north)
        {СУЩНОСТЬ --- КАК ХРАНИМ};
      \node[font=\scriptsize\bfseries,text=accent,above=0.3em] at (d.north)
        {DTO --- ЧТО ОБЕЩАЛИ};
      \draw[arra] (e) -- node[above,font=\scriptsize,text=accent]
        {\texttt{ToDto()}} node[below,font=\scriptsize,text=muted] {на границе} (d);
    \end{tikzpicture}
  \end{center}

  \slidegap
  \aside{Внутренние поля, связи, то, что нужно базе, --- в сущности. Наружу ---
    ровно нужные поля в удобной клиенту форме.}
\end{frame}

% 27
\begin{frame}[plain]
  \vote{Зачем два типа, если поля почти совпадают?}
\end{frame}

% 28
\begin{frame}{Зачем DTO}
  \slidetop
  \claim{DTO --- это \hi{контракт}, записанный типом.}

  {\renewcommand{\arraystretch}{1.6}
  \begin{tabular}{@{}L{4.6cm}L{8.2cm}@{}}
    \hi{\textbf{не отдать лишнее}}       & себестоимость, внутренние заметки, хэш
      пароля \\
    \hi{\textbf{зафиксировать контракт}} & меняется, только когда вы решили изменить
      контракт \\
    \hi{\textbf{белый список на входе}}  & клиент задаёт только то, что можно: без
      \texttt{Id}, статуса, даты \\
  \end{tabular}}

  \slidegap
  \aside{Принять сущность целиком --- заказ «оплачен», который никто не оплачивал:
    \lo{избыточная привязка}.}
\end{frame}

% 29
\begin{frame}[fragile]{Несколько DTO на ресурс}
  \vspace{0.3em}
\begin{lstlisting}[style=csharp,basicstyle=\codexs,escapechar=|]
|\mk{d1}|public sealed record OrderDto(
    int Id, string CustomerEmail, OrderStatus Status, decimal Total);
|\mk{d2}|public sealed record OrderListItemDto(int Id, OrderStatus Status);
|\mk{d3}|public sealed record CreateOrderDto(
    string CustomerEmail,
    List<OrderItemDto> Items,
    DateOnly? DeliveryDate = null);

public static class OrderMapping
{
|\mk{d4}|    public static OrderDto ToDto(this Order order) =>
        new(order.Id, order.CustomerEmail, order.Status,
            order.Items.Sum(i => i.Price * i.Quantity));
}
\end{lstlisting}
  \begin{tikzpicture}[remember picture,overlay]
    \foreach \m/\n in {d1/1,d2/2,d3/3,d4/4} \node[num] at ([xshift=11.9cm]\m) {\n};
  \end{tikzpicture}

  \slidegap
  \aside{\hi{1} ответ \quad \hi{2} строка списка \quad \hi{3} создание --- без
    \texttt{Id} и статуса \quad \hi{4} преобразование на границе, руками}
\end{frame}

% 30
\begin{frame}[plain]
  \keyline{От запроса до параметров три шага, и каждый может сломаться.\\
    Наружу уходит DTO --- и никогда сущность.}
\end{frame}

% ======================================================== 3. VALIDATION ==
\section{3. Валидация}

% 31
\begin{frame}[plain]
  \sectionopener{3}{Валидация}{Невалидный запрос --- \hi{исключение}, ответ
    формирует middleware.}
\end{frame}

% 32
\begin{frame}{Три ступени}
  \slidetop
  {\renewcommand{\arraystretch}{1.5}
  \begin{tabular}{@{}L{3.6cm}L{4.7cm}L{4.5cm}@{}}
    & \colhead{accent}{даёт} & \colhead{alt}{не даёт} \\
    \textbf{атрибуты}              & простые правила рядом с полем & правил на два поля и
      с зависимостями \\
    \textbf{свои правила}          & код: даты, суммы, условия     & ответа, куда его
      положить \\
    \hi{\textbf{FluentValidation}} & класс на DTO, правила цепочкой, ошибки по полям
      & --- \\
  \end{tabular}}

  \slidegap
  \aside{В .NET 10 атрибуты работают и в минимальном API --- одной регистрацией. Но
    ответ --- в \lo{чужом} формате.}
\end{frame}

% 33
\begin{frame}[fragile]{Валидатор}
  \vspace{0.3em}
\begin{lstlisting}[style=csharp,basicstyle=\codexs,escapechar=`]
`\mk{v1}`public sealed class CreateOrderDtoValidator
    : AbstractValidator<CreateOrderDto>
{
    public CreateOrderDtoValidator()
    {
`\mk{v2}`        RuleFor(x => x.CustomerEmail).NotEmpty().EmailAddress();
        RuleFor(x => x.Items).NotEmpty()
            .WithMessage("Order must have at least one item");
`\mk{v3}`        RuleForEach(x => x.Items).ChildRules(item =>
        {
            item.RuleFor(i => i.Quantity).GreaterThan(0);
            item.RuleFor(i => i.Price).GreaterThanOrEqualTo(0);
        });
        RuleFor(x => x.DeliveryDate)
`\mk{v4}`            .Must(d => d is null
                || d > DateOnly.FromDateTime(DateTime.UtcNow))
            .WithMessage("Delivery date must be in the future");
    }
}
\end{lstlisting}
  \begin{tikzpicture}[remember picture,overlay,every node/.style={rawlab}]
    \node at ([xshift=9.6cm]v1) {класс на один DTO};
    \node at ([xshift=9.6cm]v2) {читается как предложение};
    \node at ([xshift=9.6cm]v3) {каждая позиция};
    \node[text=accent] at ([xshift=9.6cm]v4) {свой предикат и сообщение};
  \end{tikzpicture}
\end{frame}

% 34
\begin{frame}[plain]
  \vote{Валидатор нашёл ошибки. Что делает обработчик?}
\end{frame}

% 35
\begin{frame}{Два вида плохого входа}
  \slidetop
  \versus{Ошибка контракта}{%
    \begin{cbul}{accent}
      \item клиент нарушил договор
      \item отрицательное количество
      \item исключение
      \item 422 из middleware
      \item до логики
    \end{cbul}}{Ошибка сценария}{%
    \begin{cbul}{alt}
      \item часть логики
      \item «на складе нет»
      \item возвращаемое значение
      \item решает обработчик
      \item внутри логики
    \end{cbul}}

  \slidegap
  \hrulethin{0.42}
  \rulegap
  {\large Невалидный запрос --- \hi{исключение}, а не ветка в обработчике.}
\end{frame}

% 36
\begin{frame}[fragile]{Фильтр конечной точки}
  \vspace{0.3em}
\begin{lstlisting}[style=csharp,basicstyle=\codexxs,escapechar=|]
|\mk{f1}|orders.AddEndpointFilter(async (context, next) =>
{
|\mk{f2}|    var dto = context.Arguments.OfType<CreateOrderDto>()
        .FirstOrDefault();
    if (dto is not null)
    {
|\mk{f3}|        var validator = context.HttpContext.RequestServices
            .GetRequiredService<IValidator<CreateOrderDto>>();
        var result = await validator.ValidateAsync(
            dto, context.HttpContext.RequestAborted);
        if (!result.IsValid)
|\mk{f4}|            throw new ValidationException(result.Errors);
    }
    return await next(context);
});

|\mk{f5}|// ErrorHandlingMiddleware (L8), new branches:
// ValidationException     => 422 + errors by field
// BadHttpRequestException => bad.StatusCode (400, 413, 415)
\end{lstlisting}
  \begin{tikzpicture}[remember picture,overlay,every node/.style={rawlab}]
    \node at ([xshift=9.6cm]f1) {фильтр на группе};
    \node at ([xshift=9.6cm]f2) {DTO среди аргументов};
    \node at ([xshift=9.6cm]f3) {валидатор из контейнера};
    \node[text=alt] at ([xshift=9.6cm]f4) {бросок, не возврат};
    \node[text=accent] at ([xshift=9.6cm]f5) {ветки в middleware};
  \end{tikzpicture}
\end{frame}

% 37
\begin{frame}[fragile]{Три места}
  \slidetop
  \begin{center}
    \begin{tikzpicture}[font=\footnotesize,
      tb/.style={box,minimum width=3.0cm,minimum height=1.0cm}]
      \node[tb] (t1) at (0,0)    {\textbf{валидатор}\\правила};
      \node[tb] (t2) at (3.6,0)  {\textbf{фильтр}\\проверка};
      \node[tb] (t3) at (7.2,0)  {\textbf{middleware}\\ответ};
      \node[tb,draw=accent,line width=1.0pt] (t4) at (10.8,0)
        {\textbf{обработчик}\\не знает об ошибках};
    \end{tikzpicture}
  \end{center}

\begin{lstlisting}[style=term,basicstyle=\termsm]
$ curl -s --json @bad-order.json http://localhost:5090/orders
{
  "error": "Validation failed",
  "traceId": "0HNPV3D7GA67A:00000001",
  "errors": {
    "customerEmail": [
      "'Customer Email' is not a valid email address."
    ],
    "items[0].quantity": [
      "'Quantity' must be greater than '0'."
    ]
  }
}
\end{lstlisting}

  \aside{Один формат для всех ошибок. Имена полей --- как в JSON: клиент подсветит
    \hi{нужное поле} формы.}
\end{frame}

% 38
\begin{frame}{Граница валидатора}
  \slidetop
  \versus{Валидатор}{%
    \begin{cbul}{accent}
      \item запрос сам по себе
      \item число положительное, строка не пустая
      \item без базы и внешних сервисов
      \item 422
    \end{cbul}}{Логика}{%
    \begin{cbul}{alt}
      \item запрос относительно мира
      \item «такой email уже есть»
      \item нужны данные и состояние
      \item 409 или значение
    \end{cbul}}

  \slidegap
  \hrulethin{0.42}
  \rulegap
  \aside{Размер страницы в сто тысяч и список на миллион позиций --- тоже
    \lo{невалидный} запрос: отклонить до базы. Валидация --- не авторизация.}
\end{frame}

% ======================================================== 4. JSON ==
\section{4. JSON и сериализация}

% 39
\begin{frame}[plain]
  \sectionopener{4}{JSON и сериализация}{JSON --- это строки и числа; всё
    остальное --- \hi{договорённость}.}
\end{frame}

% 40
\begin{frame}[plain]
  \vote{В C\# свойство называется \texttt{CustomerEmail}. Как оно будет называться
    в JSON, если ничего не настраивать?}
\end{frame}

% 41
\begin{frame}[fragile]{JSON и типы .NET}
  \slidetop
  {\renewcommand{\arraystretch}{1.2}\small
  \begin{tabular}{@{}L{4.0cm}L{8.8cm}@{}}
    имена                 & \hi{camelCase} при записи \\
    регистр при чтении    & не важен \\
    \texttt{null}         & пишется как \texttt{null} \\
    \texttt{enum}         & \lo{число}; строка --- по настройке \\
    даты                  & ISO 8601: \texttt{2026-11-10}, смещение, если есть \\
    \texttt{decimal}      & число с нужными знаками \\
    нет поля              & значение по умолчанию, \lo{молча} \\
  \end{tabular}}

  \slidegap
\begin{lstlisting}[style=term,basicstyle=\termsm]
$ curl -s localhost:5090/sample
{"customerEmail":"a@b.by","status":1,"date":"2026-11-10","total":12.50,"note":null}
\end{lstlisting}

  \rulegap
  \aside{Веб-умолчания --- не умолчания сериализатора: их выставили \hi{за вас}.}
\end{frame}

% 42
\begin{frame}{Числа и деньги}
  \slidetop
  \claim{JSON не различает целое и дробное. Клиент --- \lo{различает}.}

  {\renewcommand{\arraystretch}{1.6}
  \begin{tabular}{@{}L{4.4cm}L{8.4cm}@{}}
    \hi{\textbf{больше $2^{53}$}}      & клиент на JavaScript теряет младшие разряды
      --- идентификатор строкой \\
    \hi{\textbf{\texttt{decimal}}}     & уходит числом: \texttt{25.00}; клиент читает
      его в \texttt{double} \\
    \hi{\textbf{деньги}}               & строкой или целым в копейках \\
  \end{tabular}}

  \slidegap
  \aside{Решение контракта --- до первого клиента. Потом это ломающее изменение.}
\end{frame}

% 43
\begin{frame}[plain]
  \vote{Способ оплаты: карта, счёт или наличные, у каждого свои поля. Как
    записать это в JSON и прочитать обратно?}
\end{frame}

% 44
\begin{frame}[fragile]{Полиморфизм}
  \vspace{0.3em}
\begin{lstlisting}[style=csharp,basicstyle=\codexs,escapechar=|]
|\mk{y1}|[JsonPolymorphic(TypeDiscriminatorPropertyName = "kind")]
|\mk{y2}|[JsonDerivedType(typeof(CardPayment), "card")]
|\mk{y3}|[JsonDerivedType(typeof(InvoicePayment), "invoice")]
public abstract record Payment(decimal Amount);

public sealed record CardPayment(decimal Amount, string Last4)
    : Payment(Amount);
public sealed record InvoicePayment(decimal Amount, string InvoiceNumber)
    : Payment(Amount);
\end{lstlisting}
  \begin{tikzpicture}[remember picture,overlay,every node/.style={rawlab}]
    \node[text=accent] at ([xshift=9.0cm]y1) {имя дискриминатора};
    \draw[rawbrace] ([xshift=8.75cm,yshift=0.35em]y2) -- ([xshift=8.75cm,yshift=-0.3em]y3);
    \node at ([xshift=9.0cm]$(y2)!0.5!(y3)$) {значение для наследника};
  \end{tikzpicture}

  \slidegap
\begin{lstlisting}[style=json,basicstyle=\codesm]
{"kind":"card","last4":"4242","amount":100}
\end{lstlisting}

  \slidegap
  \aside{JSON не несёт типов --- тип несёт \hi{поле}, и оно становится частью
    контракта.}
\end{frame}

% 45
\begin{frame}{Почему полиморфизм болит}
  \slidetop
  {\renewcommand{\arraystretch}{1.35}\small
  \begin{tabular}{@{}L{4.4cm}L{8.4cm}@{}}
    \lo{\textbf{дискриминатор не первый}} & \lo{500}: \texttt{must specify a type
      discriminator} \\
    \lo{\textbf{неизвестное значение}}    & ошибка чтения: клиент старой версии падает
      целиком \\
    \lo{\textbf{переименовать}}           & нельзя --- никогда \\
    \lo{\textbf{документ}}                & описывает варианты неидеально \\
  \end{tabular}}

  \slidegap
  \eyebrow{Две скучные альтернативы}\par\vspace{0.5em}
  \begin{cbul}{accent}
    \item отдельные маршруты на каждый способ, у каждого плоский DTO
    \item поле-перечисление «вид» и условные правила в валидаторе
  \end{cbul}
\end{frame}

% 46
\begin{frame}[fragile]{Настройки --- один раз}
  \slidetop
  \claim{Контракт \hi{однородный}: настройки на всё приложение.}

\begin{lstlisting}[style=csharp,basicstyle=\codesm,escapechar=|]
builder.Services.ConfigureHttpJsonOptions(options =>
{
    var json = options.SerializerOptions;
|\mk{j1}|    json.Converters.Add(new JsonStringEnumConverter());
|\mk{j2}|    json.Converters.Add(new EmailJsonConverter());
|\mk{j3}|    json.RespectNullableAnnotations = true;
|\mk{j4}|    json.RespectRequiredConstructorParameters = true;
|\mk{j5}|    json.WriteIndented = builder.Environment.IsDevelopment();
});
\end{lstlisting}
  \begin{tikzpicture}[remember picture,overlay,every node/.style={rawlab}]
    \node at ([xshift=10.1cm]j1) {перечисления строками};
    \node at ([xshift=10.1cm]j2) {свой конвертер};
    \node at ([xshift=10.1cm]j3) {\texttt{null} в обязательном};
    \node at ([xshift=10.1cm]j4) {нет поля --- ошибка};
  \end{tikzpicture}

  \slidegap
  \aside{Отступы --- единственная развилка по окружению. Настройка в отдельном
    обработчике --- \lo{никогда}.}
\end{frame}

% 47
\begin{frame}[fragile]{Свой конвертер}
  \vspace{0.3em}
\begin{lstlisting}[style=csharp,basicstyle=\codexxs]
public sealed class EmailJsonConverter : JsonConverter<Email>
{
    public override Email Read(
        ref Utf8JsonReader reader, Type type, JsonSerializerOptions options) =>
        Email.TryParse(reader.GetString(), out var email)
            ? email
            : throw new JsonException("Not an email");

    public override void Write(
        Utf8JsonWriter writer, Email value, JsonSerializerOptions options) =>
        writer.WriteStringValue(value.Value);
}
\end{lstlisting}

  \rulegap
  {\renewcommand{\arraystretch}{1.3}
  \begin{tabular}{@{}L{3.4cm}l@{}}
    \lo{без конвертера} & {\codexs \{"email":\{"value":"anna@example.com"\},"name":"Anna"\}} \\
    \hi{с конвертером}  & {\codexs \{"email":"anna@example.com","name":"Anna"\}} \\
  \end{tabular}}

  \rulegap
  \aside{Тип с одним полем ведёт себя в JSON как \hi{строка}. Конвертер находится
    по типу свойства.}
\end{frame}

% 48
\begin{frame}[fragile]{Генератор кода}
  \vspace{0.3em}
\begin{lstlisting}[style=csharp,basicstyle=\codesm,escapechar=|]
|\mk{x1}|[JsonSerializable(typeof(OrderDto))]
[JsonSerializable(typeof(CreateOrderDto))]
[JsonSerializable(typeof(Payment))]
|\mk{x2}|internal sealed partial class AppJsonContext
    : JsonSerializerContext;

builder.Services.ConfigureHttpJsonOptions(options =>
|\mk{x3}|    options.SerializerOptions.TypeInfoResolverChain
        .Insert(0, AppJsonContext.Default));
\end{lstlisting}
  \begin{tikzpicture}[remember picture,overlay,every node/.style={rawlab}]
    \node at ([xshift=9.4cm]x1) {по атрибуту на тип контракта};
    \node[text=accent] at ([xshift=9.4cm]x2) {вторую половину пишет генератор};
    \node at ([xshift=9.4cm]x3) {сначала контекст, потом рефлексия};
  \end{tikzpicture}

  \slidegap
  \begin{cbul}{accent}
    \item под AOT --- \hi{обязательно}: тип, забытый в контексте, --- исключение на
      первом запросе
    \item без AOT --- быстрее старт, все типы контракта в одном месте
    \item в контексте --- DTO, а не сущности
  \end{cbul}
\end{frame}

% 49
\begin{frame}[plain]
  \statement{Сериализация --- \lo{слабое место} платформы. Кто напишет
    нормальный сериализатор --- я первым на него \hi{перейду}.}
\end{frame}

% 50
\begin{frame}[fragile]{Стриминг ответа}
  \vspace{0.3em}
\begin{lstlisting}[style=csharp,basicstyle=\codexs,escapechar=|]
|\mk{s1}|app.MapGet("/orders/export",
    (IOrderService service, CancellationToken ct) => service.ExportAsync(ct));

public async IAsyncEnumerable<OrderDto> ExportAsync(
|\mk{s2}|    [EnumeratorCancellation] CancellationToken ct)
{
    await foreach (var order in repository.ReadAllAsync(ct))
|\mk{s3}|        yield return order.ToDto();
}
\end{lstlisting}
  \begin{tikzpicture}[remember picture,overlay,every node/.style={rawlab}]
    \node at ([xshift=9.8cm]s1) {последовательность};
    \node at ([xshift=9.8cm]s2) {токен --- внутрь};
    \node[text=accent] at ([xshift=9.8cm]s3) {первый уходит сразу};
  \end{tikzpicture}

\begin{lstlisting}[style=term,basicstyle=\termsm]
$ curl -i http://localhost:5090/orders/export
HTTP/1.1 200 OK
Content-Type: application/json; charset=utf-8
Date: Tue, 10 Nov 2026 11:10:57 GMT
Server: Kestrel
Transfer-Encoding: chunked

[{"id":1,"customerEmail":"anna@example.com","status":"Paid","total":25.00},{"id":2,...
\end{lstlisting}

\end{frame}

% 51
\begin{frame}[fragile]{Цена стриминга}
  \slidetop
  \claim{Статус уходит с \lo{первым} элементом.}

\begin{lstlisting}[style=term,basicstyle=\termsm]
$ curl -s -i http://localhost:5090/orders/export
HTTP/1.1 200 OK
Content-Type: application/json; charset=utf-8
Date: Tue, 10 Nov 2026 11:10:59 GMT
Server: Kestrel
Transfer-Encoding: chunked

[{"id":1,"customerEmail":"anna@example.com","status":"Paid","total":25.00},{"id":2,...
$ echo $?
18
\end{lstlisting}

  \rulegap
  \begin{cbul}{alt}
    \item ошибка на третьем элементе --- 200 и \lo{оборванный} массив
    \item обработчик ошибок видит начатый ответ и пропускает исключение
  \end{cbul}

  \aside{Всё, что может упасть, --- \hi{до} первого элемента.}
\end{frame}

% 52
\begin{frame}[plain]
  \keyline{Сериализация --- набор решений, принятых за вас:\\
    задайте их один раз. Держите контракт плоским.}
\end{frame}

% ======================================================== 5. CONTRACT OUTSIDE ==
\section{5. Описание контракта наружу}

% 53
\begin{frame}[plain]
  \sectionopener{5}{Описание контракта наружу}{Документ \hi{генерируется} из
    кода, а не пишется рядом с ним.}
\end{frame}

% 54
\begin{frame}[fragile]{OpenAPI}
  \slidetop
  {\renewcommand{\arraystretch}{1.25}\small
  \begin{tabular}{@{}L{4.0cm}L{8.8cm}@{}}
    пути                & из шаблонов маршрутов \\
    параметры           & из типов параметров \\
    \hi{ответы}         & из \hi{типизированных результатов} \\
    схемы               & из DTO и настроек сериализатора \\
  \end{tabular}}

  \slidegap
\begin{lstlisting}[style=csharp,basicstyle=\codesm]
builder.Services.AddOpenApi();
app.MapOpenApi();   // /openapi/v1.json
\end{lstlisting}

\begin{lstlisting}[style=term,basicstyle=\termsm]
$ curl -s localhost:5090/openapi/v1.json | jq -c '.paths["/orders/{id}"].get.responses|keys'
["200","404"]
\end{lstlisting}

  \aside{Главный читатель --- генератор клиентов. Рукописная часть точна, пока её
    обновляют.}
\end{frame}

% 55
\begin{frame}[plain]
  \vote{Вы переименовали поле в DTO. Клиенты сломались? А если добавили?}
\end{frame}

% 56
\begin{frame}{Обратная совместимость}
  \slidetop
  \versus{Можно без версии}{%
    \begin{cbul}{accent}
      \item добавить поле в ответ
      \item добавить необязательный параметр
      \item добавить маршрут
      \item добавить значение перечисления*
    \end{cbul}}{Нельзя}{%
    \begin{cbul}{alt}
      \item переименовать или удалить поле
      \item изменить тип поля
      \item сделать поле запроса обязательным
      \item изменить смысл статуса или формат
    \end{cbul}}

  \slidegap
  \aside{* если клиент обязан переживать неизвестные значения --- договоритесь
    заранее.}
\end{frame}

% 57
\begin{frame}{Версия}
  \slidetop
  \claim{Версия --- это \lo{дорого}. Сначала спросите, нельзя ли добавить.}

  {\renewcommand{\arraystretch}{1.5}
  \begin{tabular}{@{}L{4.0cm}L{8.8cm}@{}}
    \hi{\textbf{в адресе}}       & префикс группы: \texttt{/api/v1}, \texttt{/api/v2} \\
    \hi{\textbf{рядом}}          & старая работает, клиенты переезжают в своём темпе \\
    \hi{\textbf{устаревание}}    & заголовки и пометки в документе; отключение ---
      в датах \\
    \lo{\textbf{цена}}           & два набора обработчиков, DTO и тестов --- годами \\
  \end{tabular}}

  \slidegap
  \aside{Контракт --- то, что вы обещали, а не то, что написали. Порядок полей и
    сортировку вы не обещали.}
\end{frame}

% 58
\begin{frame}[plain]
  \keyline{Документ выводится из кода.\\ Контракт --- обещание: добавлять можно,
    менять и удалять нельзя.}
\end{frame}

\end{document}
